% Generated draft: requires mathematical and source review.
\documentclass[11pt]{article}
\usepackage[margin=25mm]{geometry}
\usepackage{fontspec}
\setmainfont{Latin Modern Roman}
\usepackage{amsmath,amssymb,mathtools,unicode-math}
\setmathfont{Latin Modern Math}
\usepackage{enumitem,needspace,xurl,hyperref}
\setlength{\parindent}{0pt}
\setlength{\parskip}{5pt}
\setlength{\emergencystretch}{3em}
\newcommand{\N}{\mathbb{N}}
\newcommand{\Z}{\mathbb{Z}}
\newcommand{\Q}{\mathbb{Q}}
\newcommand{\R}{\mathbb{R}}
\newcommand{\C}{\mathbb{C}}
\newcommand{\F}{\mathbb{F}}
\newcommand{\A}{\mathbb{A}}
\newcommand{\PP}{\mathbb{P}}
\newcommand{\E}{\mathbb{E}}
\newcommand{\eps}{\varepsilon}
\newcommand{\dd}{\,\mathrm{d}}
\begin{document}
\section*{UnsolvedMath ID: 30003646}
\textbf{Source problem code:} catalogue entry
\subsection{Definitions and mathematical statement}
Let \(\N=\{1,2,3,\ldots\}\). Under the usual normalized convention, a function \(f:\N\to\R_{\geq0}\) is multiplicative if \(f(1)=1\) and
\[
f(mn)=f(m)f(n)\qquad\text{whenever }m,n\in\N\text{ and }\gcd(m,n)=1.
\]
For real \(x>0\), define
\[
F(x)=\sum_{\substack{n\in\N\\n\leq x}}f(n),\qquad
F_2(x)=\sum_{\substack{n\in\N\\n\leq x\\2\nmid n}}f(n).
\]
Assume that for every fixed real \(\lambda>0\),
\[
\lim_{x\to\infty}\frac{F(\lambda x)}{F(x)}=1.
\]
The limit is taken through real \(x\), and the assumption is pointwise in the fixed parameter \(\lambda\); no uniformity in \(\lambda\) is stipulated.

Does there necessarily exist \(L\in[0,1]\) such that
\[
\lim_{x\to\infty}\frac{F_2(x)}{F(x)}=L?
\]

The source does not explicitly state the normalization \(f(1)=1\). The formulation above uses the standard normalized meaning of multiplicative, which ensures \(F(x)>0\) for \(x\geq1\). If multiplicativity is instead defined solely by the coprime product identity, the identically zero function must be excluded for the displayed quotients to be defined. This convention requires confirmation in mathematical review. The supplied background mentions a predicted Euler-factor value, but supplies no formula for it; the explicit source statement asks only for existence of the limit.
\subsection{Short English statement}
For a nonnegative multiplicative function whose summatory function is slowly varying under every fixed positive dilation, must the proportion of its total weight supported on odd integers converge? The standard normalization of multiplicativity is made explicit above; the source leaves that convention implicit.
\subsection{Sources}
\begin{itemize}
\item Supplied local UnsolvedMath registry record, permanent ID \(30003646\), title ``Odd-Density Limits for Multiplicative Functions,'' source problem code \texttt{catalogue entry}; registry context: \texttt{data/unsolvedmath/problems.json} and \texttt{catalog/problems.json}. The supplied original, clean, and displayed statements have the same mathematical target.
\item Supplied report citation: \emph{Mini-Workshop: Interplay between Number Theory and Analysis for Dirichlet Series} (2018), \emph{Oberwolfach Reports}, pp.~3035--3069, DOI \texttt{10.4171/owr/2017/51}. Source URL: \url{https://doi.org/10.4171/owr/2017/51}. The supplied literature-source label is ``Oberwolfach Report 51, Hilberdink problem,'' with the same link: \url{https://doi.org/10.4171/owr/2017/51}.
\item The supplied background attributes the question to Hilberdink; the structured proposer field is empty. The supplied proposed year is 2017, whereas the report citation labels the workshop 2018. Both are retained without resolving their relationship.
\item Literature assessment checked at \texttt{2026-08-21}: the supplied curation labels the problem open and reports that the report and a subsequent search revealed no general proof or counterexample. It describes a predicted Euler-factor value without giving its formula. This is supplied dated curation, not an independently verified literature assessment or a claim that the cited source was visited for this draft.
\end{itemize}
\end{document}
